\scanpage{9}
description du groupe fondamental mixte en un point $x\in\omega(I)$
\deleted{$\unsure{\Pi_1}S_0\mathcal T_{0,I}$}
\[
\begin{aligned}
(17)\qquad\Pi_1(S_0\mathcal T_{0,I},\mathfrak S_I;x)
&=\{\alpha,l\mid\alpha\in\mathfrak S_I,\ l:x\longleftarrow\alpha x\text{ dans }S_0\mathcal T_{0,I}\}\\
&=\{(\alpha,(n_i)_{i\in I})\in\mathfrak S_I\times\mathbb Z^I\mid\text{les }n_i\text{ ont la parité définie par }\operatorname{sg}(\alpha)\}
\end{aligned}
\]
\margin{de façon générale, si $x,y\in\varepsilon\ (=\omega(I))$ les hom \unsure{dans le} $\Pi_1$ mixte de $x$ dans $y$ sont donnés par les couples $(l,\alpha)$ avec $\boxed{l:\alpha x\longrightarrow\deleted{\unsure{x}}y}$ dans le $\Pi_1$ \uline{\unsure{pur}} \unsure{(16)|(18)} -- comp. donnée (18).}
avec comme composition
\[
(18)\qquad(l,\alpha)(l',\alpha')=
\deleted{\unsure{(\alpha\alpha',}}
\bigl(\deleted{\unsure{\alpha\alpha'}}\underbrace{l\circ\alpha(l')}_{l+\alpha(l')},\deleted{\unsure{l}}\alpha\alpha'\bigr)
\]
Comme \deleted{$\unsure{\mathcal E}$} $S_0\mathcal T_{0,I}$ est déduit d'un $\mathbb Z$-groupoïde (sur $\omega(I)$ comme ens d'objets) par changement de groupe d'opérateurs interne $\mathbb Z\;-\;\mathbb Z^I$, l'extension envisagée de $\mathfrak S_I$ par $\mathbb Z^I$ peut être regardée aussi comme déduite d'une extension de $\mathfrak S_I$ par $\mathbb Z$ (savoir le groupe fond. mixte \deleted{\unsure{on a}} du \inserted{$\mathbb Z$-}groupoïde (à opérateurs $\mathfrak S_I$) précédent).

Sa description explicite est la $=$ \deleted{\unsure{que}}\inserted{que dans} (17), (18), sauf que les $l$ sont des éléments de $\mathbb Z$ au lieu de $\mathbb Z^I$ -- et l'op de $\mathfrak S_I$ sur $\mathbb Z$ \deleted{est} est triviale. Donc on voit que l'ext. envisagée n'est autre que l'extension, déduite de l'extension (14) de $\mu_2$ par $\mathbb Z$, via l'homom signature $\mathfrak S_I\;-\;\mu_2$. En résumé:

\uline{Proposition} Soit $x\in\omega(I)$. Alors $\Pi_1(S_0\mathcal T_{0,I},\mathfrak S_I;x)$ est can. isom. \unsure{à la l'extension|à l'extension} de $\mathfrak S_I$ par $\mathbb Z^I$
\mspage{8}
